\section{A Communication-Centric Roofline}
\label{sec:roofline}
\label{sec:simulator}  % keeps existing \ref{sec:simulator} links alive

\paraq{What is this section for?}
Section~\ref{sec:hybrids} showed that the fastest strategy---data parallelism, gsplit, or a
hybrid---changes with the model, the graph, the depth, the hidden size, and the interconnect.
This section introduces a \emph{roofline} on which every layer of every strategy appears as one
point, so that a single plot per setting shows \emph{what bounds each layer} and why one strategy
beats another at that layer. The roofline is a tool for identifying the \textbf{bottleneck}; it
is not a timing model, and it is not used here to choose a strategy. Section~\ref{sec:predicted}
plots it for representative settings, and Section~\ref{sec:results} checks how well its picture
agrees with the measurements.

% ---------------------------------------------------------------------------
\subsection{Purpose and composition rules}
\label{sec:cm-purpose}

\paraq{Why a roofline and not a timing model?}
GNN training time cannot be predicted exactly from first principles: kernel efficiency, launch
overheads, and contention all vary with the sampled graph. A roofline avoids this by asking a
weaker question---\emph{which resource limits this layer?}---whose answer is robust to those
details. Each resource (dense arithmetic, scattered memory reads, the interconnect) gives one
upper bound on throughput, a \emph{roof}; a layer runs no faster than the lowest roof that
applies to it, and that roof names its bottleneck.

\paraq{How are the costs of different events combined?}
The roofline rests on one rule: events that are \emph{active at the same time} overlap, so their
combined time is the \textbf{maximum} of their times (equivalently, the attainable throughput is
the \textbf{minimum} over their roofs); events that \emph{follow one another} add. The rules used
in this section are:
\begin{center}
\small
\begin{tabular}{l l l}
\toprule
Events & Treated as & Combined time \\
\midrule
gather and GEMM (one layer, one GPU) & concurrent (bottleneck) & $\max$ \\
compute and shuffle, gsplit & sequential & sum \\
compute and shuffle, pipeline & concurrent across micro-batches & $2\max$ per micro-batch \\
the $G$ GPUs of one layer & concurrent, synced by the shuffle & $\max$ over GPUs \\
successive layers & sequential & sum over layers \\
forward and backward pass & constant factor & (dropped) \\
\bottomrule
\end{tabular}
\end{center}
The first row is the standard roofline assumption; in practice the gather and GEMM kernels are
launched back to back, so the true compute time lies between the $\max$ and the sum, a factor of
at most $2$. Because of the $\max$ rules the roofline is an \emph{optimistic} bound: it places
each layer at the best throughput its bottleneck allows. This is why it identifies bottlenecks
reliably but does not predict time.

% ---------------------------------------------------------------------------
\subsection{Cost of computing one node}
\label{sec:cm-node}

\paraq{What work does computing a single node involve?}
The hidden representation $H^{l}_u$ of a destination node $u$ in layer $l$ is computed from the
hidden representations of its sampled neighbours, $H^{l-1}_v$ for $v \in N(u)$, and the layer's
weights. For GraphSAGE (mean aggregator):
\begin{align}
  H^{l}_u &= \sigma\Big( W^{\mathrm{self}}_l H^{l-1}_u
      + W^{\mathrm{neigh}}_l \tfrac{1}{|N(u)|}\textstyle\sum_{v \in N(u)} H^{l-1}_v \Big).
      \label{eq:sage}
\end{align}
For GAT:
\begin{align}
   z^{l}_v &= W_l\, H^{l-1}_v,
      \qquad v \in N(u) \cup \{u\}
      && \text{dense: one GEMM row per source node} \label{eq:gat-dense}\\
   \alpha_{uv} &= \underset{v \in N(u)}{\mathrm{softmax}}\Big(
         \mathrm{LeakyReLU}\big(a_l^{\top}\,[\, z^{l}_u \,\|\, z^{l}_v \,]\big)\Big)
      && \text{sparse: one scalar score per edge} \label{eq:gat-attn}\\
   H^{l}_u &= \sigma\Big( \sum_{v \in N(u)} \alpha_{uv}\, z^{l}_v \Big)
      && \text{sparse: weighted gather over } N(u) \label{eq:gat-agg}
\end{align}
Both models interleave two kinds of work. \emph{Dense} work is vector--matrix products
(a row times $W$) and vector--vector products (the attention dot product), over contiguous
memory; it is bound by the arithmetic rate. \emph{Sparse} work is the aggregation over $N(u)$,
which reads neighbour rows at scattered addresses; it is bound by the bandwidth of scattered
DRAM reads.

\paraq{Which constants describe a model?}
Every hidden layer has width $h$, so a hidden row is $s_H = h\,b$ bytes with $b$ bytes per value
($b{=}4$ for fp32). Let $e$ be the number of sampled edges into $u$ (its fanout). A model is
described by five constants, read off its layer equations (Table~\ref{tab:cm-constants}):
\begin{itemize}\itemsep1pt
  \item $\alpha$: hidden rows read per edge by the sparse kernels;
  \item $\beta_e, \beta_n$: vector--matrix products ($h^2$ multiply--adds each) per edge and per node;
  \item $\gamma_e, \gamma_n$: vector--vector products ($h$ multiply--adds each) per edge and per node.
\end{itemize}
The constants depend only on the model; the graph enters only through $e$.

\paraq{What are the dense and sparse work of one node?}
Counting a multiply--add as two FLOPs, the dense work and the sparse (gather) traffic of $u$ are
\begin{align}
  F(e) &= 2\big[(\beta_e\, e + \beta_n)\, h^{2} + (\gamma_e\, e + \gamma_n)\, h\big]
       \quad\text{FLOPs}, \label{eq:dense-cost}\\
  S(e) &= \alpha\, e\, s_H \quad\text{bytes.} \label{eq:sparse-cost}
\end{align}
With the measured dense rate $\hat{P}$ and gather bandwidth $\hat{B}_{\mathrm{mem}}$
(Table~\ref{tab:cm-microbench}) and the bottleneck rule of Section~\ref{sec:cm-purpose}, the
compute time of one node is
\begin{equation}
  T^{\mathrm{cmp}}_u \;=\; \max\!\Big( \frac{F(e)}{\hat{P}},\; \frac{S(e)}{\hat{B}_{\mathrm{mem}}} \Big).
  \label{eq:tcmp}
\end{equation}

\paraq{Why measured rates rather than peak rates?}
GNN kernels rarely approach the peak FLOP rate or the peak DRAM bandwidth, and how close they get
depends on the model, $h$, and the sparsity. Peak rates would understate compute time and make
recomputation look cheaper than it is. We therefore use rates measured by microbenchmarks that
run the same kernels on sampled blocks, per hidden size.

\paraq{What about the backward pass?}
The constants describe the forward pass. We assume the backward pass multiplies the compute
\emph{and} the shuffled bytes (gradients travel the reverse path) by about the same factor, so it
moves every point by the same amount and does not change which roof binds. It is left out.

\begin{table}[h]
\centering
\caption{Model constants of Eqs.~\eqref{eq:dense-cost}--\eqref{eq:sparse-cost} (forward pass,
per destination node). \placeholder{} = to be derived from the layer equations.}
\label{tab:cm-constants}
\small
\begin{tabular}{l c c c c c}
\toprule
Model & $\alpha$ & $\beta_e$ & $\beta_n$ & $\gamma_e$ & $\gamma_n$ \\
\midrule
SAGE & \placeholder & \placeholder & \placeholder & \placeholder & \placeholder \\
GAT  & \placeholder & \placeholder & \placeholder & \placeholder & \placeholder \\
\bottomrule
\end{tabular}
\end{table}

\begin{table}[h]
\centering
\caption{Measured hardware rates (microbenchmarks on one A100 node; all-to-all over 4 GPUs).
\placeholder{} = to be measured.}
\label{tab:cm-microbench}
\small
\begin{tabular}{l l c c c c c}
\toprule
Rate & Microbenchmark & $h{=}256$ & $512$ & $1024$ & $2048$ & $4096$ \\
\midrule
$\hat{P}$ (TFLOP/s) & one layer's GEMM on a sampled block
  & \placeholder & \placeholder & \placeholder & \placeholder & \placeholder \\
$\hat{B}_{\mathrm{mem}}$ (GB/s) & neighbour gather, fanout $e$
  & \placeholder & \placeholder & \placeholder & \placeholder & \placeholder \\
$\hat{B}_{\mathrm{net}}$ (GB/s), PCIe-emu & NCCL all-to-all of $s_H$-byte rows
  & \placeholder & \placeholder & \placeholder & \placeholder & \placeholder \\
$\hat{B}_{\mathrm{net}}$ (GB/s), NVLink & NCCL all-to-all of $s_H$-byte rows
  & \placeholder & \placeholder & \placeholder & \placeholder & \placeholder \\
\bottomrule
\end{tabular}
\end{table}

% ---------------------------------------------------------------------------
\subsection{The axes: why bandwidth is the reference}
\label{sec:cm-axes}

\paraq{Why not the classical roofline?}
A classical roofline plots attained FLOP/s against arithmetic intensity, and its flat roof is the
peak FLOP rate. For GNN layers that roof is not a fixed line: the achievable FLOP rate depends on
the model and on the sparsity of the sampled block. The interconnect bandwidth, in contrast, is a
property of the machine alone. We therefore measure throughput in \emph{bytes of hidden state},
the unit the interconnect moves, so that the interconnect becomes the flat, uniform roof and the
model-dependent compute appears as sloped roofs.

\paraq{What is on the $y$-axis?}
The \textbf{useful throughput per GPU}: bytes of hidden state per second that a GPU contributes
to the layer's output. ``Useful'' counts each distinct output row once, so a row computed
redundantly by several GPUs earns credit only once. Let $u$ be the number of distinct destination
rows the layer must produce across all $G$ GPUs; each GPU's useful share is
$U = (u/G)\,s_H$ bytes, and
\begin{equation}
  y \;=\; \frac{U}{T} \quad \text{bytes/s per GPU},
  \label{eq:cm-y}
\end{equation}
where $T$ is the layer time.

\paraq{What is on the $x$-axis?}
The \textbf{compute per useful byte}: FLOPs a GPU spends per byte of useful output,
\begin{equation}
  x \;=\; \frac{n\, F(\bar{e})}{U} \quad \text{FLOP/byte},
  \label{eq:cm-x}
\end{equation}
where $n$ is the number of destination rows the GPU actually computes and $\bar{e}$ their average
fanout. If every row is computed once and the GPUs are balanced, $n = u/G$ and
$x = F(\bar{e})/s_H$; redundancy ($n > u/G$) and imbalance both move the point to the right.

\paraq{Why not FLOPs per \emph{communicated} byte?}
That is the classical arithmetic intensity with the network in place of DRAM. Data parallelism
communicates nothing, so its intensity would be infinite and its points would leave the plot.
Normalising by useful bytes keeps every strategy on the same axes.

\paraq{How do the model and the hidden size move a point?}
From Eq.~\eqref{eq:dense-cost}, for balanced, non-redundant work,
$x = 2\big[(\beta_e\bar{e}+\beta_n)\,h + (\gamma_e\bar{e}+\gamma_n)\big]/b$, which grows linearly in
$h$. A model with more dense work per row (larger $\beta$) sits further right, and doubling $h$
roughly doubles $x$.

% ---------------------------------------------------------------------------
\subsection{The three roofs}
\label{sec:cm-roofs}

\paraq{What is the dense roof?}
If the GEMMs were the only cost, a GPU would spend $nF/\hat{P}$ seconds producing $U$ useful bytes,
so by Eqs.~\eqref{eq:cm-y}--\eqref{eq:cm-x}
\begin{equation}
  y \;\le\; \frac{U\hat{P}}{nF} \;=\; \frac{\hat{P}}{x}.
  \label{eq:roof-dense}
\end{equation}
On log--log axes this is a line of slope $-1$.

\paraq{What is the gather roof?}
If the gather were the only cost, the GPU would spend $nS/\hat{B}_{\mathrm{mem}}$ seconds. Define
the \emph{gather ratio}
\begin{equation}
  g \;=\; \frac{S(\bar{e})}{F(\bar{e})} \;=\; \frac{\alpha\,\bar{e}\,s_H}{F(\bar{e})}
  \quad\text{bytes gathered per FLOP},
  \label{eq:gather-ratio}
\end{equation}
which depends only on the model, $h$, and $\bar{e}$. The gathered bytes per useful byte are then
$nS/U = g\,x$, and
\begin{equation}
  y \;\le\; \frac{U\hat{B}_{\mathrm{mem}}}{nS} \;=\; \frac{\hat{B}_{\mathrm{mem}}}{g\,x}.
  \label{eq:roof-gather}
\end{equation}
This is also a line of slope $-1$, \emph{parallel} to the dense roof: gather traffic and FLOPs
both scale with the number of rows computed, so the two roofs never cross. Whichever is lower
binds everywhere for that (model, $h$, $\bar{e}$).

\paraq{Which compute roof binds?}
By the $\max$ rule of Eq.~\eqref{eq:tcmp}, the compute roof is the lower of the two:
\begin{equation}
  R^{\mathrm{cmp}}(x) \;=\; \frac{\hat{P}_{\mathrm{bind}}}{x},
  \qquad
  \hat{P}_{\mathrm{bind}} \;=\; \min\!\Big(\hat{P},\; \frac{\hat{B}_{\mathrm{mem}}}{g}\Big).
  \label{eq:roof-cmp}
\end{equation}
The layer is \textbf{FLOP-bound} when $g < \hat{B}_{\mathrm{mem}}/\hat{P}$ and
\textbf{gather-bound} otherwise. Since $F \propto h^2$ and $S \propto h$, $g \propto 1/h$:
small hidden sizes tend to be gather-bound and large ones FLOP-bound. $\hat{P}_{\mathrm{bind}}$
is the achievable FLOP rate of the layer---this is the precise sense in which ``different models
and sparsities have different achievable FLOP rates''.

\paraq{What is the interconnect roof?}
Pulling one hidden row over the interconnect delivers $s_H$ bytes in $s_H/\hat{B}_{\mathrm{net}}$
seconds, so a layer that pulls one row per useful row is bounded by the flat line
\begin{equation}
  y \;\le\; \hat{B}_{\mathrm{net}}.
  \label{eq:roof-net}
\end{equation}
There is one such line per interconnect (emulated PCIe and NVLink), identical for every model.
A layer that receives $m$ rows on the GPU has a \emph{pull ratio}
\begin{equation}
  r \;=\; \frac{m}{u/G} \quad\text{rows received per useful row},
  \label{eq:pull-ratio}
\end{equation}
and its own interconnect ceiling is $\hat{B}_{\mathrm{net}}/r$: the flat roof shifted down by the
partition's cut (for $r>1$) or up (for $r<1$).

\paraq{Where is the ridge?}
For a layer with pull ratio $r$, the compute and transfer times are equal where
$R^{\mathrm{cmp}}(x) = \hat{B}_{\mathrm{net}}/r$, i.e.\ at
\begin{equation}
  x^{*}(r) \;=\; \frac{r\,\hat{P}_{\mathrm{bind}}}{\hat{B}_{\mathrm{net}}}.
  \label{eq:ridge}
\end{equation}
For $r{=}1$ this is the crossing of the sloped compute roof with the flat $\hat{B}_{\mathrm{net}}$
line. The ridge on NVLink lies far to the left of the ridge on emulated PCIe, because
$\hat{B}_{\mathrm{net}}$ is much larger.

\begin{figure}[h]
  \centering
  \fbox{\parbox{0.85\linewidth}{\centering\small
  \expdiag{schematic roofline. $x$: compute per useful byte (log); $y$: useful bytes/s per GPU
  (log). Dashed: dense roof $\hat{P}/x$ and gather roof $\hat{B}_{\mathrm{mem}}/(gx)$ (parallel);
  solid: the binding compute roof $\hat{P}_{\mathrm{bind}}/x$. Flat: $\hat{B}_{\mathrm{net}}$ for
  PCIe-emu and NVLink; dotted: $\hat{B}_{\mathrm{net}}/r$ for one layer. Ridge $x^*$ marked.
  Regions of Section~\ref{sec:cm-regions} shaded. One gsplit, one pipeline and one data-parallel
  point of the same layer, as in Section~\ref{sec:cm-strategies}.}\\[4pt]
  \placeholder}}
  \caption{Schematic communication-centric roofline (placeholder).}
  \label{fig:cm-schematic}
\end{figure}

% ---------------------------------------------------------------------------
\subsection{A layer as a point}
\label{sec:cm-point}

\paraq{Which quantities does sampling provide?}
A strategy fixes which rows each GPU computes and which rows it receives; one sampling pass
reveals both, without features or training. For strategy $s$, layer $\ell$ (counted from the
seeds, as in Section~\ref{sec:hybrid-dp}), and GPU $q$, sampling gives
\begin{itemize}\itemsep1pt
  \item $n^{s}_{\ell,q}$: destination rows computed;
  \item $e^{s}_{\ell,q}$: edges aggregated, so $\bar{e}^{s}_{\ell,q} = e^{s}_{\ell,q}/n^{s}_{\ell,q}$;
  \item $m^{s}_{\ell,q}$: rows received over the interconnect ($m = 0$ for data parallelism);
  \item $u^{s}_{\ell}$: distinct destination rows of the layer over all GPUs.
\end{itemize}

\paraq{How long does the layer take on each GPU?}
Applying the composition rules of Section~\ref{sec:cm-purpose},
\begin{align}
  T^{\mathrm{cmp}}_{\ell,q} &= \max\!\Big(\frac{n_{\ell,q} F(\bar{e}_{\ell,q})}{\hat{P}},\;
       \frac{\alpha\, e_{\ell,q}\, s_H}{\hat{B}_{\mathrm{mem}}}\Big), \qquad
  T^{\mathrm{net}}_{\ell,q} = \frac{m_{\ell,q}\, s_H}{\hat{B}_{\mathrm{net}}},
  \label{eq:cm-gpu-parts}\\
  T^{s}_{\ell,q} &= \begin{cases}
     T^{\mathrm{cmp}}_{\ell,q} & \text{data parallel},\\
     T^{\mathrm{cmp}}_{\ell,q} + T^{\mathrm{net}}_{\ell,q} & \text{gsplit (split layer)},\\
     \max\!\big(T^{\mathrm{cmp}}_{\ell,q},\, T^{\mathrm{net}}_{\ell,q}\big) & \text{pipeline (split layer)}.
  \end{cases}
  \label{eq:cm-gpu-time}
\end{align}
The pipeline case follows from its two half-size micro-batches: each micro-batch has half the
compute and half the shuffle, and one computes while the other shuffles, so
\[
  T^{\mathrm{pipe}} = 2\max\!\big(T^{\mathrm{cmp}}/2,\; T^{\mathrm{net}}/2\big)
  = \max\!\big(T^{\mathrm{cmp}}, T^{\mathrm{net}}\big).
\]

\paraq{Which GPU represents the layer?}
The shuffle synchronises the GPUs, so the layer ends when its slowest GPU does. The point is
taken from the \textbf{busiest GPU},
$q^{*} = \arg\max_q T^{s}_{\ell,q}$, which keeps load imbalance visible instead of averaging it
away. The layer's point is
\begin{equation}
  x^{s}_\ell = \frac{n_{\ell,q^*}\, F(\bar{e}_{\ell,q^*})}{(u^{s}_\ell/G)\, s_H},
  \qquad
  y^{s}_\ell = \frac{(u^{s}_\ell/G)\, s_H}{T^{s}_{\ell,q^*}},
  \qquad
  r^{s}_\ell = \frac{m_{\ell,q^*}}{u^{s}_\ell/G}.
  \label{eq:cm-point}
\end{equation}
An $L$-layer model gives $L$ points per strategy.

\paraq{What else does sampling give for free?}
The activation memory of the busiest GPU, about $\sum_\ell n_{\ell,q^*} s_H$ times a
model-dependent factor, so configurations that cannot fit (cf.\ the out-of-memory cells of
Section~\ref{sec:hybrid-empirical}) can be marked on the plot.

% ---------------------------------------------------------------------------
\subsection{Where each strategy lands}
\label{sec:cm-strategies}

\paraq{Where does a data-parallel layer sit?}
It has no transfer, so $T = T^{\mathrm{cmp}}$ and, by Eqs.~\eqref{eq:roof-cmp}
and~\eqref{eq:cm-point}, $y = \hat{P}_{\mathrm{bind}}/x$ exactly: the point lies \emph{on} the
binding compute roof. Data parallelism pays through its $x$: each GPU computes
$n = \delta\, u/G$ rows, where
\begin{equation}
  \delta_\ell \;=\; \frac{G\, n_{\ell,q^*}}{u_\ell} \;\ge\; 1
  \label{eq:cm-redundancy}
\end{equation}
combines redundancy (rows computed by several GPUs) with imbalance. The point slides $\delta$
times further right along the roof, and therefore $\delta$ times lower, than a non-redundant
layer of the same model. Since the trees of different GPUs overlap more as they grow, $\delta$
rises with distance from the seeds, and the outer data-parallel layers drift right.

\paraq{Where does a gsplit layer sit?}
Each row is computed once, so $x$ is near its non-redundant value, but compute and shuffle add:
\begin{equation}
  \frac{1}{y} \;=\; \frac{x}{\hat{P}_{\mathrm{bind}}} \;+\; \frac{r}{\hat{B}_{\mathrm{net}}}.
  \label{eq:cm-gsplit-point}
\end{equation}
The point lies \emph{below} both the compute roof and its interconnect ceiling
$\hat{B}_{\mathrm{net}}/r$, and at least half as high as the lower of the two. It is furthest
below (by a factor of $2$) at the ridge, where the two terms are equal.

\paraq{Where does a pipeline layer sit?}
With $T = \max(T^{\mathrm{cmp}}, T^{\mathrm{net}})$,
\begin{equation}
  y \;=\; \min\!\Big( \frac{\hat{P}_{\mathrm{bind}}}{x},\; \frac{\hat{B}_{\mathrm{net}}}{r} \Big),
  \label{eq:cm-pipe-point}
\end{equation}
so the point lies \emph{on} the lower roof, directly above the gsplit point of the same layer.
The vertical gap is the pipeline's gain,
\begin{equation}
  \frac{y^{\mathrm{pipe}}}{y^{\mathrm{split}}} \;=\; \frac{T^{\mathrm{cmp}} + T^{\mathrm{net}}}{\max(T^{\mathrm{cmp}}, T^{\mathrm{net}})}
  \;=\; 1 + \frac{\min(T^{\mathrm{cmp}}, T^{\mathrm{net}})}{\max(T^{\mathrm{cmp}}, T^{\mathrm{net}})}
  \;\in\; [1, 2],
  \label{eq:cm-pipe-gain}
\end{equation}
which is largest, $2$, when $T^{\mathrm{cmp}} = T^{\mathrm{net}}$, i.e.\ at the ridge
$x^{*}(r)$. Far from the ridge one term dominates and the extra overlap has little to hide.
The pipeline is thus best for layers near the ridge.

\paraq{Where does a data-parallel hybrid sit?}
With split point $k$, layers $\ell < k$ are data-parallel points (on the compute roof, shifted
right by $\delta_\ell$) and layers $\ell \ge k$ are gsplit points (below both roofs). On one plot
the hybrid's $L$ points are therefore a mixture of the two series, and at each layer the higher of
the data-parallel and gsplit points is the faster strategy for that layer.

\paraq{How are two strategies compared at one layer?}
By height. Both points of a layer refer to (approximately) the same useful output $U$, so the
higher point has the shorter layer time. Data parallelism wins a layer when sliding right by
$\delta$ along the compute roof costs less height than the transfer term of
Eq.~\eqref{eq:cm-gsplit-point}. Comparing heights assumes $u^{\mathrm{dp}}_\ell \approx
u^{\mathrm{split}}_\ell$; the two samplings are not identical, and the difference is listed among
the assumptions below.

% ---------------------------------------------------------------------------
\subsection{Regions of the roofline}
\label{sec:cm-regions}

\paraq{Which regions does the roofline divide into?}
The roofs and the ridge split the plane into regions, each naming a bottleneck and the strategy
that the bottleneck favours. For a layer with pull ratio $r$ and band parameter
$\beta_{\mathrm{band}} > 1$ (\placeholder):
\begin{center}
\footnotesize
\begin{tabular}{l l l l}
\toprule
Region & Condition & Bottleneck & Favours \\
\midrule
I.\ Communication-bound & $x < x^{*}(r)/\beta_{\mathrm{band}}$
   & interconnect & DP / larger $k$ \\
II.\ Balanced (ridge band) & $x^{*}(r)/\beta_{\mathrm{band}} \le x \le \beta_{\mathrm{band}}\, x^{*}(r)$
   & neither dominates & pipeline \\
III.a Compute-bound, gather & $x > \beta_{\mathrm{band}}\, x^{*}(r)$, $g \ge \hat{B}_{\mathrm{mem}}/\hat{P}$
   & DRAM gather & gsplit \\
III.b Compute-bound, FLOP & $x > \beta_{\mathrm{band}}\, x^{*}(r)$, $g < \hat{B}_{\mathrm{mem}}/\hat{P}$
   & dense arithmetic & gsplit \\
\bottomrule
\end{tabular}
\end{center}

\paraq{Why does each region favour that strategy?}
\begin{itemize}\itemsep2pt
  \item \textbf{Region I (left of the ridge).} The compute roof is high, so computing is cheap per
        useful byte, and the transfer term of Eq.~\eqref{eq:cm-gsplit-point} dominates the gsplit
        point. Recomputing instead of pulling costs a slide of $\delta$ along a roof that is far
        above the interconnect ceiling, so data parallelism (or a larger split point $k$) wins
        unless $\delta$ is large enough to push the data-parallel point past the ridge.
  \item \textbf{Region II (around the ridge).} Compute and transfer are comparable. gsplit loses up
        to half its throughput to the sum of the two (Eq.~\eqref{eq:cm-gsplit-point}), and the
        pipeline recovers up to $2\times$ (Eq.~\eqref{eq:cm-pipe-gain}).
  \item \textbf{Region III (right of the ridge).} The layer is compute-bound; the transfer is small
        next to the compute. Redundant computation is expensive here because sliding right by
        $\delta$ lowers the point in proportion, so gsplit, which computes each row once, wins.
        The sub-region (gather- or FLOP-bound) says which compute resource a faster kernel would
        have to improve; it does not change the strategy.
\end{itemize}

\paraq{How do the settings of Section~\ref{sec:hybrids} move a layer between regions?}
\begin{itemize}\itemsep1pt
  \item \textbf{Interconnect.} A slower interconnect lowers $\hat{B}_{\mathrm{net}}$ and moves the
        ridge $x^{*}$ to the right (Eq.~\eqref{eq:ridge}), so the same layer moves from Region III
        on NVLink towards Region I on emulated PCIe.
  \item \textbf{Hidden size.} Larger $h$ moves the point right ($x \propto h$) and lowers $g$
        (towards FLOP-bound); with $\hat{P}_{\mathrm{bind}}$ fixed this moves the layer towards
        Region III.
  \item \textbf{Model.} More dense work per row (GAT) moves the point right of a SAGE layer of the
        same $h$.
  \item \textbf{Depth / hop.} Outer layers have larger frontiers and more overlap between GPUs:
        $\delta_\ell$ grows (data-parallel points move right) and $r_\ell$ changes with the
        partition's cut, which moves the ridge. Different layers of one model can therefore fall
        in different regions---the condition under which a per-layer hybrid can beat both pure
        strategies.
\end{itemize}

% ---------------------------------------------------------------------------
\subsection{Assumptions to check}
\label{sec:cm-assumptions}

\paraq{Which assumptions does the roofline rest on?}
\begin{enumerate}\itemsep1pt
  \item Gather and GEMM of a layer overlap ($\max$); in practice they run back to back, so true
        compute time is up to $2\times$ the modelled time.
  \item The backward pass scales compute and shuffled bytes by the same factor.
  \item All hidden layers have width $h$; input-feature loading and the first layer's feature
        width are not modelled.
  \item The per-edge dense products $\beta_e e h^2$ are charged to every destination node,
        i.e.\ no sharing of source-row products between destinations (an over-count for GAT,
        whose $z_v$ is computed once per source row).
  \item $u^{\mathrm{dp}}_\ell \approx u^{\mathrm{split}}_\ell$ when comparing heights at one
        layer.
  \item The pipeline has exactly two half-size micro-batches; fill and drain, and the lower
        efficiency of half-size GEMMs, are ignored.
  \item Sampling and feature loading are outside the plot; they add the same per-epoch cost to
        every strategy only to first order.
  \item The microbenchmark rates of Table~\ref{tab:cm-microbench} are independent of the layer
        (a single $\hat{P}$, $\hat{B}_{\mathrm{mem}}$, $\hat{B}_{\mathrm{net}}$ per $h$).
\end{enumerate}

\paraq{Conclusion of this section.}
Each layer of each strategy becomes one point: its $x$ is the compute it spends per useful byte,
its $y$ the useful bytes it delivers per second on its busiest GPU. The interconnect is a flat
roof shared by all models; dense and gather work are parallel sloped roofs whose lower one is the
achievable compute rate. Data-parallel layers sit on the compute roof, shifted right by
redundancy; gsplit layers sit below both roofs; pipeline layers sit on the lower roof. The ridge
divides the plane into a communication-bound region, a balanced band, and a compute-bound
region, and the region a layer falls in names its bottleneck. Section~\ref{sec:predicted} draws
this roofline for representative settings.
